\section{Interacción Natural y Experiencia de Usuario (UX)}

\subsection{Integración de voz y texto}
El reconocimiento de voz se apoya en la API estándar del navegador
(\textit{Speech Recognition}). El micrófono se activa pulsando un 
botón en la interfaz o manteniendo presionada la barra espaciadora. 
Para dar \textit{feedback} inmediato al usuario y evitar la 
sensación de latencia, se utilizan los resultados intermedios 
(\textit{interim results}); de este modo, el texto va apareciendo 
en la pantalla a medida que se habla, enviándose al servidor justo 
cuando el usuario termina la locución.

Para la lectura de las respuestas se utiliza la API de síntesis de 
voz (\textit{Speech Synthesis}). No se espera a que el 
modelo genere todo el párrafo, sino que la emisión de audio se 
sincroniza con el flujo de datos (\textit{streaming}) y se va 
reproduciendo por frases completas, utilizando los signos de puntuación como 
delimitadores. Esto soluciona problemas frecuentes de 
lectura con cifras o decimales y evita cortes a 
mitad de palabra. Además, el usuario puede personalizar el tono 
(\textit{pitch}), la velocidad (\textit{rate}) y la voz del 
asistente desde un panel de configuración, guardándose sus 
preferencias localmente en el navegador (\textit{localStorage}).

\subsection{Generación de respuestas en streaming}

Toda la comunicación en tiempo real entre el \textit{cliente} y 
el \textit{servidor} se realiza a través de una conexión WebSocket 
abierta de forma continua. Gracias a este canal, el usuario no 
tiene que esperar a que el servidor procese la respuesta completa 
para empezar a leerla, sino que el texto va apareciendo en la interfaz
frase a frase, a medida que el modelo las completa. La unidad no es el
\textit{token} sino la frase, porque es lo que la síntesis de voz puede
leer sin cortar palabras ni tropezar con los decimales.

Además del texto de la respuesta, pueden transmitirse 
actualizaciones relacionadas con los datos mostrados en pantalla,
como cambios en el saldo o instrucciones para representar tablas 
y gráficos.

Esta estrategia reduce el tiempo de espera percibido y permite que la interfaz
comience a mostrar información mientras el procesamiento de la consulta
continúa en segundo plano. De esta forma, la generación de la 
respuesta y la actualización de los distintos elementos de 
la interfaz pueden realizarse de forma progresiva.

La Figura~\ref{fig:turno} recoge la secuencia completa de un turno en el que
el usuario hace una consulta sobre su histórico. El texto llega frase a frase y
la síntesis de voz empieza con la primera; la tabla o el gráfico se generan en
un tramo posterior, cuando la respuesta ya está completa.

\begin{figure}[htbp]
\centering
\begin{tikzpicture}[
    font=\small,
    actor/.style={draw=#1, fill=#1!7, rounded corners=2pt, align=center,
                  minimum width=2.8cm, minimum height=0.95cm, inner sep=3pt},
    actor/.default=subsectionblue,
    vida/.style={draw=gray!45, dashed},
    ida/.style={-{Stealth[length=2mm]}, thick, draw=gray!70!black},
    vuelta/.style={-{Stealth[length=2mm]}, thick, dashed, draw=gray!70!black},
    msg/.style={font=\scriptsize, above, inner sep=2pt},
    paso/.style={draw=gray!60, fill=white, rounded corners=2pt, font=\scriptsize,
                 align=center, inner sep=3pt},
    nota/.style={font=\scriptsize\itshape, text=ugrred, align=right},
]

% Columnas
\def\N{0}     % navegador
\def\B{4.2}   % backend
\def\L{8.4}   % modelo
\def\H{12.6}  % herramientas y BD

% Tramo posterior a la respuesta
\fill[subsectionblue!6, rounded corners=3pt] (-1.5,-9.5) rectangle (14.1,-13.65);

% Líneas de vida
\foreach \x in {\N,\B,\L,\H} \draw[vida] (\x,-0.5) -- (\x,-13.65);

\node[font=\scriptsize\bfseries, text=subsectionblue, anchor=north east, align=right,
      fill=subsectionblue!6, inner sep=2pt]
     at (14.0,-9.6) {Después de \texttt{fin\_respuesta}:\\el usuario ya está oyendo la respuesta};

% Cabeceras
\node[actor]        at (\N,0) {Navegador};
\node[actor]        at (\B,0) {Backend\\[-1pt]\scriptsize\texttt{agent.py}};
\node[actor=ugrred] at (\L,0) {Modelo\\[-1pt]\scriptsize\texttt{qwen3:8b}};
\node[actor=gray]   at (\H,0) {Herramientas\\[-1pt]\scriptsize y base de datos};

% --- Respuesta ------------------------------------------------------------
\draw[ida]    (\N,-1.0) -- node[msg] {pregunta}                          (\B,-1.0);
\node[paso]   at (\B,-1.8) {atajos deterministas (saldo, Bizum): no aplican};
\draw[ida]    (\B,-2.6) -- node[msg] {prompt + historial + herramientas} (\L,-2.6);
\draw[vuelta] (\L,-3.4) -- node[msg] {pide \texttt{consultar\_movimientos}} (\B,-3.4);
\draw[ida]    (\B,-4.2) -- node[msg] {\texttt{tool\_inicio}, \texttt{sql}} (\N,-4.2);
\draw[ida]    (\B,-5.0) -- node[msg] {SQL en solo lectura}               (\H,-5.0);
\draw[vuelta] (\H,-5.8) -- node[msg] {filas}                             (\B,-5.8);
\draw[ida]    (\B,-6.6) -- node[msg] {resultado de la herramienta}       (\L,-6.6);
\draw[vuelta] (\L,-7.4) -- node[msg] {respuesta en \textit{streaming}}   (\B,-7.4);
\draw[ida]    (\B,-8.2) -- node[msg] {\texttt{texto}, una frase por evento} (\N,-8.2);
\node[nota, anchor=east] at (\N-0.1,-8.2) {la voz\\empieza};
\draw[ida]    (\B,-9.0) -- node[msg] {\texttt{fin\_respuesta}}           (\N,-9.0);

% --- Motor visual ---------------------------------------------------------
\node[paso]   at (\B,-10.1) {¿hay algo que mostrar?\\regla determinista: sí};
\draw[ida]    (\B,-11.0) -- node[msg] {\texttt{visual\_generando}}        (\N,-11.0);
\draw[ida]    (\B,-11.8) -- node[msg] {llamada dedicada (3 filas)} (\L,-11.8);
\draw[vuelta] (\L,-12.5) -- node[msg] {tipo, spec sin datos y motivo}    (\B,-12.5);
\draw[ida]    (\B,-13.25) -- node[msg] {\texttt{grafico} con datos} (\N,-13.25);

\end{tikzpicture}
\caption{Secuencia de un turno con consulta al histórico y gráfico. La voz
empieza con la primera frase del texto. El motor visual solo se ejecuta cuando
la respuesta ya está completa, de modo que la segunda llamada al modelo no
retrasa lo que el usuario oye.}
\label{fig:turno}
\end{figure}


\subsection{Interfaz adaptativa y accesible}

La interfaz no se limita a mostrar texto plano. Cuando una consulta 
devuelve datos comparables, el sistema lo detecta y enriquece la 
respuesta con gráficos o tablas interactivas.

El diseño original delegaba en el modelo de lenguaje la decisión 
de cuándo pintar estos elementos visuales (controlado mediante 
la variable \texttt{VISUALES\_AL\_LLM}). Sin embargo, las métricas 
revelaron que darle esta responsabilidad al LLM bloqueaba la 
emisión de la voz: el modelo solía detenerse a generar la compleja 
estructura de datos del gráfico antes de empezar a redactar el 
texto en español, lo que disparaba la latencia inicial hasta unos 12,4 segundos de media.

Para resolverlo, la configuración actual oculta por defecto las 
herramientas visuales al modelo principal. Ahora es el 
\textit{backend} quien evalúa de forma determinista el volumen y 
tipo de los datos obtenidos en la consulta para decidir si estos 
son graficables. Si lo son, lanza una llamada dedicada al LLM para 
construir la especificación del gráfico, pero lo hace estrictamente
después de emitir el evento \texttt{fin\_respuesta}. Dado que 
la memoria de la tarjeta gráfica empleada no permite paralelizar 
dos peticiones al LLM, esta segunda solicitud visual debe 
esperar su turno en cola obligatoriamente.

La ventaja obtenida es que el usuario escucha la respuesta hablada 
de inmediato (el tiempo hasta la voz se redujo a unos 6,7 segundos) y, 
mientras atiende a la explicación del asistente, la interfaz 
muestra un indicador de carga hasta que el gráfico o la tabla 
aparecen en pantalla. 